% =========================================================
% PAGE 2: EXECUTIVE SUMMARY — ARABIC (START)
% =========================================================

% 1. Adjust Geometry for RTL Sidebar
\newgeometry{right=4.5cm, left=2cm, top=1cm, bottom=2.5cm}

% ---------------------------------------------------------
% TITLE & INTRO
% ---------------------------------------------------------
\noindent
{\fontsize{24}{28}\selectfont \textbf{\color{ecestitle} ملخص تنفيذي}}
\vspace{0.8em}
\label{sec:ar_exec_summary}

\noindent
{\fontsize{11}{15}\selectfont \color{textdark}
\ECESContent{exec_intro_ar}{[TODO: أدخل الفقرة التمهيدية التي تصف نطاق مسح مقياس الأعمال للعدد \IssueNumber. اذكر الربع \CurrentQuarterText، والربع القادم \NextQuarterText، والربع السابق \PreviousQuarterText، والربع المناظر \CorrespQuarterText.]}
}

% ---------------------------------------------------------
% SUBHEADER AND ARROW LEGEND
% ---------------------------------------------------------
\noindent
\begin{minipage}[c]{0.55\textwidth}
    \raggedright
    \fontsize{12}{14}\selectfont
    \textbf{\color{ecesteal} \underline{تقييم الأداء والتوقعات وفقا للمؤشر الإجمالي}}
\end{minipage}%
\hfill
\begin{minipage}[c]{0.40\textwidth}
    \centering
    
\includegraphics[width=\linewidth]{arrow_ar.png}
\end{minipage}

\vspace{0.8em}

% ---------------------------------------------------------
% STYLED HIGHLIGHTS BOX
% ---------------------------------------------------------
\noindent
\begin{tikzpicture}
    \setRTL
    \node[
        draw=ecesteal,
        line width=1.5pt,
        fill=ecesteal!5,
        rounded corners=8pt,
        inner sep=12pt,
        align=justify,
        text width=12cm
    ] (box) {
        \fontsize{11}{15}\selectfont \color{textdark}

        \ECESContent{exec_box_overall_ar}{[TODO: أدخل نقطتين أو ثلاثاً عن مؤشر أداء الأعمال الإجمالي ومؤشر التوقعات لربع \CurrentQuarterText.]}

        \vspace{0.8em}

        \textbf{\color{textblue} \underline{وفقا للحجم:}} \\
        \ECESContent{exec_box_size_ar}{[TODO: أدخل نقطة أو نقطتين عن أداء المؤشر وفقاً لحجم الشركات (كبيرة / ص و م) لربع \CurrentQuarterText.]}

        \textbf{\color{textpurple} \underline{وفقا للقطاع:}} \\
         \ECESContent{exec_box_sector_ar}{[TODO: أدخل نقطة أو نقطتين عن أداء المؤشر وفقاً للقطاعات لربع \CurrentQuarterText.]}

        \textbf{\color{textgreen} \underline{التحديات:}} \\
        \ECESContent{exec_box_challenges_ar}{[TODO: أدخل نقطة أو نقطتين عن أبرز معوقات الأعمال لربع \CurrentQuarterText.]}
    };
\end{tikzpicture}

\vspace{0.8em}

% ---------------------------------------------------------
% BPI ANALYSIS TEXT AND CHART
% ---------------------------------------------------------
%% TODO: Upload ch1.png (Arabic) — مؤشر أداء الأعمال الإجمالي via Chart Manager
\noindent
\begin{minipage}[t]{0.55\textwidth}
    \vspace{0pt}
    {\fontsize{11}{15}\selectfont \color{textdark}
    \ECESContent{exec_bpi_analysis_ar}{[TODO: أدخل تحليل مؤشر أداء الأعمال لربع \CurrentQuarterText مع المقارنة بالربع السابق \PreviousQuarterText والربع المناظر \CorrespQuarterText.]}
    }
\end{minipage}%
\hfill
\begin{minipage}[t]{0.40\textwidth}
    \vspace{0pt}
    \centering
    \includegraphics[width=\linewidth, keepaspectratio]{ch1.png}
\end{minipage}

\vspace{0.8em}

\noindent
{\fontsize{11}{15}\selectfont \color{textdark}
\ECESContent{exec_expectations_ar}{[TODO: أدخل تحليل مؤشر توقعات الأداء للربع القادم \NextQuarterText.]}
}

\clearpage

% =========================================================
% PAGE 3: DETAILED ANALYSIS (SIZE & SECTOR) — ARABIC
% =========================================================

%% TODO: Upload ch2.png (Arabic) — مؤشر الأداء وفقاً للحجم via Chart Manager
\noindent
\begin{minipage}[t]{0.55\textwidth}
    \vspace{0pt}
    {\fontsize{11}{15}\selectfont \textbf{\color{textblue} \underline{وفقا للحجم:}}}
    \vspace{0.6em}

    {\fontsize{11}{15}\selectfont \color{textdark}
    \ECESContent{exec_size_detail_ar}{[TODO: أدخل تحليل تفصيلي لأداء المؤشر وفقاً لحجم الشركات لربع \CurrentQuarterText.]}
    }
\end{minipage}%
\hfill
\begin{minipage}[t]{0.40\textwidth}
    \vspace{0pt}
    \centering
    \includegraphics[width=\linewidth, keepaspectratio]{ch2.png}
\end{minipage}

\vspace{1.5em}

\noindent
{\fontsize{11}{15}\selectfont \textbf{\color{textpurple} \underline{وفقا للقطاع:}}} \vspace{0.6em}

\noindent
{\fontsize{11}{15}\selectfont \color{textdark}
\ECESContent{exec_sector_summary_ar}{[TODO: أدخل ملخص أداء المؤشر وفقاً للقطاعات لربع \CurrentQuarterText. اذكر أفضل وأسوأ القطاعات أداءً.]}
}

\vspace{1.2em}

%% TODO: Upload ch3.png (Arabic) — مؤشر الأداء وفقاً للقطاعات via Chart Manager
\noindent
\includegraphics[width=\linewidth, keepaspectratio]{ch3.png}

\vspace{1.2em}

\noindent
{\fontsize{11}{15}\selectfont \color{textdark}
\ECESContent{exec_manufacturing_ar}{[TODO: أدخل تحليل أداء قطاع الصناعات التحويلية لربع \CurrentQuarterText.]}

\vspace{0.8em}

\ECESContent{exec_best_sector_ar}{[TODO: أدخل تحليل أفضل قطاع أداءً لربع \CurrentQuarterText.]}
}

\clearpage

% =========================================================
% PAGE 4: CHALLENGES & MACROECONOMIC DEVELOPMENTS — ARABIC
% =========================================================

\noindent
{\fontsize{11}{15}\selectfont \textbf{\color{textgreen} \underline{التحديات والأولويات من منظور مجتمع الأعمال:}}} \vspace{0.8em}

\noindent
{\fontsize{11}{15}\selectfont \textbf{\color{textdark}
\ECESContent{exec_challenges_headline_ar}{[TODO: أدخل أبرز معوق واجه كافة الشركات خلال ربع \CurrentQuarterText.]}
}} \vspace{0.8em}

\noindent
{\fontsize{11}{15}\selectfont \color{textdark}
\ECESContent{exec_challenges_body_ar}{[TODO: أدخل قائمة بأبرز 4–5 معوقات تواجه الأعمال خلال \CurrentQuarterText مع شرح موجز لكل منها.]}

\vspace{0.8em}

\ECESContent{exec_sme_comparison_ar}{[TODO: أدخل المقارنة بين معوقات الشركات الصغيرة والمتوسطة والكبيرة خلال \CurrentQuarterText.]}

\vspace{0.8em}

\ECESContent{exec_priorities_ar}{[TODO: أدخل أبرز الأولويات لتحسين بيئة الأعمال من منظور العينة خلال \CurrentQuarterText.]}
}

\vspace{2em}

\noindent
{\fontsize{11}{15}\selectfont \textbf{\color{black} \underline{أبرز التطورات الاقتصادية الكلية: *}}} \vspace{0.5em}

\begin{itemize} \itemsep0.6em
    \item \fontsize{11}{15}\selectfont \textbf{على الصعيد الدولي:} \ECESContent{exec_macro_global_ar}{[TODO: أدخل ملخصاً موجزاً للأوضاع الاقتصادية العالمية ذات الصلة بربع \CurrentQuarterText.]}

    \item \fontsize{11}{15}\selectfont \textbf{على الصعيد المحلي:} \ECESContent{exec_macro_local_ar}{[TODO: أدخل ملخصاً موجزاً للأوضاع الاقتصادية الكلية في مصر خلال ربع \CurrentQuarterText.]}
\end{itemize}

\vfill

\noindent
{\fontsize{9}{11}\selectfont
* لمزيد من التفاصيل حول أحدث المستجدات المالية العالمية والإقليمية والمحلية، يُرجى الرجوع إلى التقرير المالي الشهري (\textbf{\color{eceslink} \underline{لقطة أسواق المال}}).
}

\restoregeometry
